\chapter{Containerization with Docker and Docker Compose}\label{ch:docker}

I did all the Docker work in this chapter on the project's Ubuntu 24.04 VM
(\texttt{vm-23L-0764-A2-Ops}, Docker Engine 29.9, Compose v5.7), connected over SSH from Azure Cloud
Shell. Every screenshot shows the \texttt{azureuser@vm-23L-0764-A2-Ops} prompt with the date and
time. The same files also work with Docker Desktop on Windows without changes, and GitHub Actions
builds and smoke-tests both images on every push (\cref{sec:6-6}).

\section{Docker warm-up}
Before putting my own project in containers I practised with the ready-made
\texttt{nginx:1.27-alpine} image (\file{labs/docker_warmup.md}). The command
\cmd{docker run -d --name warmup -p 8081:80 -v .../labs/warmup-site:/usr/share/nginx/html:ro nginx:1.27-alpine}
runs it in the background (\cmd{-d}), maps port 8081 on the host to port 80 in the container
(\cmd{-p}) and mounts a folder with my own \file{index.html} read-only (\cmd{-v ...:ro}). The page
then shows up at \texttt{http://localhost:8081}. \cmd{docker logs} shows an nginx access-log line
for each refresh, \cmd{docker exec -it warmup sh} opens a shell inside the running container and
\cmd{docker ps -a} lists the containers. \cmd{docker stop} and \cmd{docker rm} remove the
container, and \cmd{docker images} and \cmd{docker rmi} list and remove the image.

\shot{s6_1_run}{Run with port mapping and volume; page served in the browser.}
\shot{s6_1_logs_exec}{Logs and a shell inside the container.}
\shot{s6_1_cleanup}{Listing and removing containers and images.}

\begin{description}[leftmargin=2.6cm,style=nextline]
\item[Image] a read-only template built from a Dockerfile: a file system plus metadata (command,
environment, exposed ports). An image does not run by itself.
\item[Container] a running or stopped instance of an image. It is an isolated process with its own
view of the file system, its own network and process space, and a thin writable layer that is lost
when the container is removed.
\item[Layer] each Dockerfile instruction that changes the file system creates one read-only layer,
and an image is a stack of them. Layers are cached and shared between images, so an instruction
that has not changed is not rebuilt.
\item[Volume] storage kept outside the container's writable layer, so it survives when the
container is removed. A bind mount maps a folder of the host instead.
\item[Port mapping] forwards a port on the host to a port inside the container's private network
(\texttt{host:container}). Without it the service cannot be reached from outside.
\end{description}

\section{Dockerfile for the web app}\label{sec:6-2}
The \href{\GitHubURL/blob/main/Dockerfile}{\file{Dockerfile}} has its instructions in this order:
\begin{enumerate}
\item \texttt{FROM python:3.11-slim}, a slim Debian without compilers or docs (about 130\,MB,
compared with about 1\,GB for \texttt{python:3.11}). All my dependencies come as wheels, so I do not
need a compiler.
\item \texttt{ENV} settings: no \texttt{.pyc} files, unbuffered output so logs appear in
\cmd{docker logs} straight away, and \texttt{PIP\_NO\_CACHE\_DIR=1} so the pip cache does not end
up inside a layer.
\item \texttt{COPY requirements-*.txt} and then \texttt{RUN pip install}, before any source code.
Docker reuses a cached layer as long as the instruction and the files it copies have not changed.
The requirements rarely change and the code changes all the time, so with this order editing a
template or a route reuses the cached TensorFlow layer (about 1\,GB) and the rebuild takes seconds.
If the source were copied first, every edit would invalidate the cache and reinstall everything.
\item \texttt{useradd vg} (uid 10001) and later \texttt{USER vg}, so the app does not run as root.
A compromised process then cannot change the image's files or install software.
\item \texttt{COPY src app} copies the code, owned by \texttt{vg}. The model is not copied here;
it is mounted at run time (\cref{sec:6-3}).
\item \texttt{HEALTHCHECK}: every 30\,s a small Python call requests \texttt{/health} (the slim
image has no \texttt{curl}), and Docker marks the container unhealthy after 3 failures.
\item \texttt{CMD gunicorn ...} starts the production WSGI server. The exec form makes gunicorn
PID 1, so it receives \texttt{SIGTERM} and shuts down cleanly.
\end{enumerate}
\file{.dockerignore} leaves \file{.git}, \file{.env}, \file{data/}, \file{mlruns/},
\file{outputs/}, the notebooks, the tests and the report out of the build context. The upload to
the Docker daemon stays small, and secrets or the dataset cannot end up in a layer. Secrets are
only passed in at run time, as environment variables from the untracked \file{.env}, and
\cmd{docker history} of the image shows none of them.

I install \texttt{tensorflow-cpu==2.20.0}. The default \texttt{tensorflow} wheel on Linux comes
with the CUDA and cuDNN libraries (several GB), which a CPU-only server never uses, and a single
MobileNetV2 prediction only takes tens of milliseconds on a CPU anyway. A build argument,
\texttt{WITH\_TF=false}, builds the web image without TensorFlow at all for
\texttt{INFERENCE\_MODE=remote}, where the model runs in the cloud API.

\shot{s6_2_whoami}{The process inside the container runs as the unprivileged user \texttt{vg} (uid 10001).}

My first container started gunicorn with \texttt{--preload}, which loads the app in the master
process before forking the workers so that start-up errors show up immediately. The health checks
passed, but the first prediction hung until gunicorn's timeout killed the worker. TensorFlow starts
its thread pools when it initialises, and \texttt{fork()} only copies the thread that calls it, so
the worker inherited locks held by threads that no longer existed and deadlocked. Without
\texttt{--preload} the worker loads the model itself after the fork, and it is still loaded only
once per process.

\section{Compose stack}\label{sec:6-3}
\href{\GitHubURL/blob/main/docker-compose.yml}{\file{docker-compose.yml}} works as follows:
\begin{itemize}
\item There are two services, \texttt{web} and \texttt{db} (\texttt{postgres:16-alpine}), on a user
network called \texttt{backend}. Compose's internal DNS resolves service names, so the web app
connects to \texttt{postgresql+psycopg2://\ldots@db:5432/visionguard}, where \texttt{db} is the
service name and not an IP address.
\item A named volume, \texttt{pgdata}, holds the database files, so the history survives
\cmd{docker compose down} and image upgrades. Only \cmd{down -v} deletes it.
\item \texttt{db} has a health check (\cmd{pg_isready}) and \texttt{web} declares
\texttt{depends\_on: db: condition: service\_healthy}, so the app only starts once PostgreSQL
accepts connections. A plain \texttt{depends\_on} only waits for the container to start, which is
too early. The app also retries its first connection.
\item The database port is not published. \texttt{db} has no \texttt{ports:} entry, so only
containers on \texttt{backend} can reach it, and it is invisible from my machine and the network.
\item Passwords come from the untracked \file{.env}. \texttt{\$\{POSTGRES\_PASSWORD:?\}} makes
Compose refuse to start if the variable is missing, and \file{.env.example} is committed with
placeholders.
\end{itemize}

The model reaches the web container through a mount: \file{./models} is mounted read-only at
\file{/app/models}. Promoting a new champion then only takes
\cmd{python scripts/export_champion.py} and \cmd{docker compose restart web}, with no rebuild, and
the 9.6\,MB model is not copied into every version of the image. Because of \texttt{:ro} the app
cannot change it. For the cloud API I did the opposite and built the model into the image
(\cref{ch:deploy}), because that image is shipped to another machine and should be self-contained
and never change after it is built.

\shot{s6_3_ps_healthy}{\cmd{docker compose ps}: both services healthy.}
\shot{s6_3_port_closed}{The database port is closed from outside (\cmd{docker compose port db 5432} is empty; a connection to localhost:5432 is refused).}

\section{Image analysis}
\begin{table}[H]\centering\small
\begin{tabular}{l r}
\toprule
Image & Size \\
\midrule
\texttt{visionguard-web:naive} (\file{docker/Dockerfile.naive}) & 6.06\,GB (1.87\,GB compressed content) \\
\texttt{visionguard-web:latest} (optimised) & 2.16\,GB (422\,MB) \\
\texttt{visionguard-web:remote} (\texttt{WITH\_TF=false}) & 385\,MB (90.8\,MB) \\
\bottomrule
\end{tabular}
\caption{Image sizes from \cmd{docker images}.}\end{table}
The naive image uses the full \texttt{python:3.11} base and the GPU \texttt{tensorflow} wheel with
its CUDA libraries, keeps the pip cache and runs as root. \cmd{docker history} shows that the
biggest layer by far is the \texttt{pip install} of TensorFlow and the serving libraries
(1.59\,GB), followed by the base image. My own code is only about 150\,KB (\file{app} 82\,KB,
\file{src} 66\,KB), and the user and metadata layers are a few KB each. For the optimised image I
switched to the slim base and \texttt{tensorflow-cpu}, turned off the pip cache and added the
ignore file.

\shot{s6_4_sizes}{\cmd{docker images}: before and after.}
\shot{s6_4_history}{\cmd{docker history}: the TensorFlow layer dominates.}

Together these changes took the web image from 6.06\,GB down to 2.16\,GB, a 64\% reduction. Almost
all of that comes from the CUDA libraries that the GPU wheel pulls in. Building without TensorFlow
for \texttt{INFERENCE\_MODE=remote}, where the model runs in the cloud API, brings it down to
385\,MB, 94\% smaller than the naive image.

\section{Operations and troubleshooting}
For day-to-day operations I used \cmd{docker compose stop web}, \cmd{start}, \cmd{restart db},
\cmd{docker compose logs --tail 30 web}, \cmd{docker compose exec web sh} and
\cmd{docker compose exec db psql -U visionguard}.
\shot{s6_5_ops}{Start/stop/restart, logs and shells in both services.}

To practise troubleshooting I broke the stack on purpose. I changed \texttt{POSTGRES\_PASSWORD} in
\file{.env} and recreated only the web container, and \cmd{docker compose ps} showed it restarting.
Using only the logs, \cmd{docker compose logs web} showed my retry loop
(\texttt{database not ready (OperationalError), retry n/10}) followed by
\texttt{FATAL: password authentication failed for user "visionguard"}. So the network and the
database were working, since the database answered, and the credentials were wrong. One detail
surprised me: the PostgreSQL image only applies \texttt{POSTGRES\_PASSWORD} when the volume is
initialised for the first time, so the database still had the old password and only the client
side had changed. To fix it I restored the value in \file{.env} and ran
\cmd{docker compose up -d --force-recreate web}. The service was healthy again and the history was
still there.
\shot{s6_5_failure_logs}{Log excerpt of the failure.}
\shot{s6_5_fixed}{After the fix: both services healthy.}

\section{Tool lab: GitHub Actions (CI)}\label{sec:6-6}
For the toy part I used a separate orphan branch of the repository, \texttt{ci-toy}, which shares
no history or files with \texttt{main} (its content is \file{labs/ci_toy/}). A workflow runs on
every push with a matrix of Python 3.11 and 3.12, so the same job runs twice in parallel, once per
version. I changed an assertion to a wrong value, which gave a red run, and then fixed it, which
gave a green one.
\shot{s6_6_toy_red}{Toy workflow: failing test, red run.}
\shot{s6_6_toy_green}{Toy workflow: fixed, green run on both Python versions.}

For VisionGuard, \file{.github/workflows/ci.yml} runs three jobs on every push:
\begin{enumerate}
\item Lint and tests on Python 3.11 and 3.12: \cmd{ruff check .} and then \cmd{pytest}. The 27
tests cover upload validation, all routes and the DB trail with a fake predictor, the inference
module, and the API's key check with a tiny model.
\item Tests against PostgreSQL, using a \texttt{postgres:16} service container, so the trail is
tested on the same database engine as in production.
\item A Docker job that builds the web and API images (with layer caching), starts them, waits for
\texttt{/health}, and checks that the process user is \texttt{vg} and that the API returns 401
without a key. It uses the committed champion model, and only generates a tiny dummy model if that
file is missing.
\end{enumerate}
The badge in \file{README.md} shows the status of the latest run.

Secrets such as a DagsHub token or a deploy key belong in the repository's Settings, under
Secrets and variables, then Actions. GitHub stores them encrypted, the workflow refers to them as
\texttt{\$\{\{ secrets.NAME \}\}}, and GitHub injects them at run time and masks them in the logs.
If I wrote them into the workflow file, anyone who can read the repository would see them, they
would stay in the Git history forever, and every fork would copy them.

\shot{s6_6_project_green}{VisionGuard CI: all jobs green.}
\shot[0.7]{s6_6_badge}{CI status badge in the README.}
